% Self-contained research section; standard amsmath/amsthm packages suffice.
% Bibliography keys and source notes are supplied in centered_cubic_notes.txt.
% No claim of historical priority is made for the lowest cubic moment.

\section{All regular odd centered cubic trigamma moments}
\label{cc:sec}

The canonical centered-square chapter evaluates all regular nonpositive
even moments of $\psi'(n+1)^2$.  Its cubic counterpart has the opposite
parity and a slightly larger, but still fixed, coordinate space.  We give
an explicit formula for every regular negative odd value, ordinary
absolutely convergent identities realizing these values, and a finite
polylogarithmic expression for their exponential generating function.

The cubic series itself has classical antecedents.  In particular,
Furdui evaluates $\sum_{n\ge1}n\psi'(n)^3$ in
\cite[Theorem 1(b)]{cc:furdui2015}; Hoffman's sums-of-tails calculus
provides a general multiple-zeta framework \cite{cc:hoffman2016}.
The contribution here is the entire centered polynomial-subtracted
family, its finite Bernoulli formula, and its explicit generator.
The lowest convergent evaluation is included as its normalization,
without a historical-priority claim.

\subsection{Regular values and ordinary convergent sums}

Put
\begin{equation}\label{cc:definitions}
 A_n=\psi'(n+1)=\zeta(2)-H_n^{(2)},\qquad
 x_n=n+\tfrac12,\qquad
 \mathcal U(s)=\sum_{n=0}^{\infty}A_n^3x_n^{-s}.
\end{equation}
The last series converges normally in $\Re s>-2$.  Define rational
numbers
\begin{equation}\label{cc:g}
 g_k=\sum_{\substack{a+b+c=k\\a,b,c\ge0}}
 B_{2a}(1/2)B_{2b}(1/2)B_{2c}(1/2),\qquad k\ge0.
\end{equation}
Thus $g_0=1$, $g_1=-1/4$, and $g_2=13/120$.
The Bernoulli convention throughout is $B_1=-1/2$.

\begin{proposition}[Continuation and the correct parity]
\label{cc:continuation}
The function $\mathcal U$ has a meromorphic continuation to the complex
plane, with simple poles at $s=-2-2k$, $k\ge0$, and residues $g_k$.
Every $s=-(2r+1)$, $r\ge0$, is regular.  Writing
\begin{equation}\label{cc:Mr}
 U_r=\mathcal U(-2r-1),
\end{equation}
one has the ordinary convergent identity
\begin{equation}\label{cc:ordinary}
 U_r=\sum_{n=0}^{\infty}
 \left\{x_n^{2r+1}A_n^3-
       \sum_{k=0}^{r-1}g_kx_n^{2r-2-2k}\right\}.
\end{equation}
The inner sum is empty for $r=0$; every bracket is $O(x_n^{-2})$.
\end{proposition}

\begin{proof}
The centered Hurwitz expansion gives, to arbitrary finite order,
\[
 A_n\sim\sum_{j\ge0}B_{2j}(1/2)x_n^{-2j-1},\qquad
 A_n^3\sim\sum_{k\ge0}g_kx_n^{-2k-3}.
\]
Subtracting finitely many terms and adding their Hurwitz-zeta sums
continues $\mathcal U$ to successive left half-planes.  The only
possible poles are the stated simple poles, with the stated residues.
For $j\ge1$, $B_{2j}(1/2)$ has sign $(-1)^j$, so every summand
in $g_k$ has sign $(-1)^k$; thus $g_k\ne0$.
At $s=-2r-1$ the subtracted sums are
$g_k\zeta(2k+2-2r,1/2)=0$ for $0\le k<r$.
Indeed $\zeta(-2m,1/2)=0$ for all $m\ge0$.
This proves \eqref{cc:ordinary} and its convergence.
\end{proof}

\subsection{A finite formula in four constant coordinates}

Set
\begin{align}
 C_1&=9\zeta(3)-\frac{17}{4}\zeta(4),
       \label{cc:C1}\\
 R_j&=-\frac{3}{2(2j-1)}
       \sum_{h=0}^{j-2}(2h+1)\binom{2j-1}{2h+3}
          B_{2j-4-2h}B_{2h}
       +\frac14\delta_{j,3},\qquad j\ge2,
       \label{cc:Rj}\\
 C_j&=9B_{2j-2}\zeta(3)-\frac12\delta_{j,2}\zeta(2)+R_j,
       \qquad j\ge2.
       \label{cc:Cj}
\end{align}
Here $\delta$ is the Kronecker delta.  The first few coordinates are
\begin{align}
 C_2&=\frac32\zeta(3)-\frac12\zeta(2)-\frac12,
       \notag\\
 C_3&=-\frac3{10}\zeta(3)-\frac25,
       \notag\\
 C_4&=\frac3{14}\zeta(3)-\frac5{56},
       \notag\\
 C_5&=-\frac3{10}\zeta(3)+\frac7{45}.
       \label{cc:Cexamples}
\end{align}

\begin{theorem}[All centered cubic moments]
\label{cc:all}
For every integer $r\ge0$,
\begin{equation}\label{cc:allformula}
 \boxed{\displaystyle
 U_r=\frac1{2r+2}\sum_{j=1}^{r+1}
      \binom{2r+2}{2j}B_{2r+2-2j}(1/2)C_j.}
\end{equation}
In particular every ordinary sum in \eqref{cc:ordinary} belongs to
\[
 \mathbb Q+\mathbb Q\zeta(2)+\mathbb Q\zeta(3)+\mathbb Q\zeta(4).
\]
No multiple-zeta coordinate of increasing weight is required as $r$
increases.
\end{theorem}

For example,
\begin{align}
 U_0&=\frac92\zeta(3)-\frac{17}{8}\zeta(4),
       \label{cc:U0}\\
 U_1&=-\frac18\zeta(2)-\frac34\zeta(3)
          +\frac{17}{32}\zeta(4)-\frac18,
       \label{cc:U1}\\
 U_2&=\frac5{48}\zeta(2)+\frac{47}{160}\zeta(3)
          -\frac{119}{384}\zeta(4)+\frac3{80},
       \label{cc:U2}\\
 U_3&=-\frac{49}{384}\zeta(2)-\frac{257}{1120}\zeta(3)
          +\frac{527}{1536}\zeta(4)-\frac{99}{4480}.
       \label{cc:U3}
\end{align}
Thus \eqref{cc:U1} evaluates
$\sum_{n\ge0}(x_n^3\psi'(n+1)^3-1)$, and
\eqref{cc:U2} evaluates
$\sum_{n\ge0}(x_n^5\psi'(n+1)^3-x_n^2+1/4)$.
These are convergent series; the displayed subtractions are part of
the identities.

\subsection{Proof by finite summation and explicit cancellation}

Let $\operatorname{CT}_N$ denote the coefficient of
$N^0(\log N)^0$ in a cutoff asymptotic expansion.  We use it only
after finite identities and standard asymptotic expansions.
Every partial sum of a nonnegative even power of $x_n$ has zero
constant coefficient as a polynomial in $N$.  Consequently
\eqref{cc:ordinary} implies
\begin{equation}\label{cc:cutoff}
 U_r=\operatorname{CT}_N\sum_{n=0}^{N-1}x_n^{2r+1}A_n^3.
\end{equation}
The centered power sum is the even polynomial
\begin{equation}\label{cc:centeredsum}
 \sum_{n=0}^{k-1}x_n^{2r+1}
 =\frac{B_{2r+2}(k+1/2)-B_{2r+2}(1/2)}{2r+2}
 =\frac1{2r+2}\sum_{j=1}^{r+1}
      \binom{2r+2}{2j}B_{2r+2-2j}(1/2)k^{2j}.
\end{equation}
Since $A_{k-1}=A_k+k^{-2}$, finite summation by parts reduces
\eqref{cc:allformula} to
\begin{equation}\label{cc:Ccutoff}
 C_j=\operatorname{CT}_N\left\{
 3\sum_{k=1}^Nk^{2j-2}A_k^2
 +3\sum_{k=1}^Nk^{2j-4}A_k
 +\sum_{k=1}^Nk^{2j-6}+N^{2j}A_N^3\right\}.
\end{equation}
It remains to prove that this expression is
\eqref{cc:C1}--\eqref{cc:Cj}.

\paragraph{The first two coordinates.}
The elementary tail rearrangements
\[
 \sum_{k=1}^{\infty}A_k^2=3\zeta(3)-\frac52\zeta(4),\qquad
 \sum_{k=1}^{\infty}\frac{A_k}{k^2}=\zeta_{\mathrm{MZ}}(2,2)=\frac34\zeta(4)
\]
give \eqref{cc:C1}.  For the first equality, include $k=0$ and
partition the two tail indices to obtain
$2\zeta_{\mathrm{MZ}}(2,1)+\zeta(3)=3\zeta(3)$, then remove $A_0^2$.

For $j=2$, write
$U_m(N)=\sum_{k=1}^Nk^mA_k$ and
$V_m(N)=\sum_{k=1}^Nk^mA_k^2$.
The same finite summation identities give
\begin{align*}
 \operatorname{CT}_N U_0(N)&=\gamma-\zeta(2)+1,\\
 \operatorname{CT}_N V_2(N)&=-\gamma+\frac12\zeta(2)
                         +\frac12\zeta(3)-\frac23,\\
 \operatorname{CT}_N(N^4A_N^3)&=-\frac32.
\end{align*}
For clarity, the first line follows from
$U_0(N)=H_N-H_N^{(2)}+NA_N$.
The second follows by using
$\sum_{k<\ell}k^2=\ell^3/3-\ell^2/2+\ell/6$ in the
quadratic summation identity below; the needed additional constants are
$\operatorname{CT}_NU_1(N)=1/4-\gamma/2$ and
$U_{-1}(\infty)=\zeta(3)$.
Substitution in \eqref{cc:Ccutoff} yields the displayed $C_2$.

\paragraph{The general coordinate.}
Let $m=2j-2\ge4$ and put
\[
 F_m(N)=\sum_{k=1}^Nk^m
 =\sum_{\substack{a=1\\a\text{ odd}}}^{m+1}f_aN^a+\frac12N^m,
 \qquad
 f_a=\frac1{m+1}\binom{m+1}{a}B_{m+1-a}.
\]
Then
\[
 F_m(k-1)=\sum_{\substack{a=1\\a\text{ odd}}}^{m+1}f_ak^a-\frac12k^m.
\]
Finite summation by parts gives the exact relation
\begin{equation}\label{cc:quadraticcutoff}
 V_m(N)=F_m(N)A_N^2+
 2\sum_{k=1}^N\frac{F_m(k-1)A_k}{k^2}
 +\sum_{k=1}^N\frac{F_m(k-1)}{k^4}.
\end{equation}
Its term $-U_{m-2}(N)$ cancels exactly against the corresponding
linear-tail term in \eqref{cc:Ccutoff}.  This cancellation occurs
before constants or divergent harmonic terms are taken.

For every positive odd $q$, another finite summation gives
\begin{equation}\label{cc:oddlinear}
 \operatorname{CT}_NU_q(N)
   =\frac{2-q}{4}B_{q-1}-\frac12\delta_{q,1}\gamma.
\end{equation}
Indeed, if $F_q(N)=\sum_{k=1}^Nk^q$, then
\[
 U_q(N)=F_q(N)A_N+\sum_{k=1}^N\frac{F_q(k-1)}{k^2}.
\]
Use the uncentered expansion
$A_N\sim\sum_{a\ge1}B_{a-1}N^{-a}$.
Only the $N^2$ and $N^q$ terms of $F_q$ can contribute to the
constant when $q\ge3$, yielding $(2-q)B_{q-1}/4$.
The case $q=1$ follows directly from
$F_1(N)=(N^2+N)/2$ and $(N-H_N)/2$.

In \eqref{cc:quadraticcutoff}, the $a=1$ terms contribute
$9f_1\zeta(3)=9B_m\zeta(3)$ to \eqref{cc:Ccutoff}.
The $a=3$ harmonic constants cancel by \eqref{cc:oddlinear}; all
remaining nonnegative-power sums have constant zero.
For odd $a\ge3$,
\[
 [N^{-a}]A_N^2=-B_{a-3}.
\]
Combining these boundary coefficients with \eqref{cc:oddlinear}
leaves the rational sum
\[
 \frac32\sum_{\substack{a=3\\a\text{ odd}}}^{m+1}
          f_a(2-a)B_{a-3},
\]
apart from the even boundary coefficient.  The latter is exactly
\begin{equation}\label{cc:boundaryexception}
 \frac32[N^{-m}]A_N^2+[N^{-m-2}]A_N^3
       =\frac14\delta_{m,4}.
\end{equation}
To see this, all odd-index Bernoulli numbers except $B_1=-1/2$
vanish.  The terms containing exactly one $B_1$ cancel the ordinary
even-even convolution; the only surviving possibility is the
$B_1^2/B_1^3$ contribution at $m=4$.
Substitute $a=2h+3$ and $m=2j-2$ to obtain
\eqref{cc:Rj}--\eqref{cc:Cj}.  This proves
Theorem~\ref{cc:all}.

\subsection{A fixed coordinate space and a rational cubic identity}

Treat $\zeta(2),\zeta(3),\zeta(4),1$ as formal coordinates.
The coefficient matrix of $U_0,U_1,U_2,U_3$ in this order is
\begin{equation}\label{cc:matrix}
 M=\begin{pmatrix}
 0&9/2&-17/8&0\\
 -1/8&-3/4&17/32&-1/8\\
 5/48&47/160&-119/384&3/80\\
 -49/384&-257/1120&527/1536&-99/4480
 \end{pmatrix},\qquad
 \det M=\frac{51}{81920}.
\end{equation}
Thus these first four values generate the entire formal coordinate
space of the cubic family, and every later value is an explicit
rational combination of them.  This is a coefficient-space statement;
it asserts no unproved arithmetic independence of the four constants.

One particularly simple elimination is
\begin{equation}\label{cc:rationalmoment}
 11U_0+198U_1+512U_2+224U_3=-\frac{21}{2}.
\end{equation}
Consequently, with
$P(x)=11x+198x^3+512x^5+224x^7$, one obtains the entirely rational
ordinary sum
\begin{equation}\label{cc:rationalsum}
 \boxed{\displaystyle
 \sum_{n=0}^{\infty}\left\{
 P(x_n)\psi'(n+1)^3-224x_n^4-456x_n^2-\frac{1414}{15}
 \right\}=-\frac{21}{2}.}
\end{equation}
Its summand is $O(x_n^{-2})$.  Equation \eqref{cc:rationalsum}
follows from the four proved values and their prescribed subtractions,
so no numerical relation search or independence assumption enters it.

Another elimination isolates the odd zeta coordinate:
\begin{equation}\label{cc:zeta3moment}
37U_0+2136U_1+11344U_2+7168U_3=252\zeta(3).
\end{equation}
With $P_3(x)=37x+2136x^3+11344x^5+7168x^7$, this becomes
\begin{equation}\label{cc:zeta3sum}
 \boxed{\begin{aligned}
 \sum_{n=0}^{\infty}\Bigl\{P_3(x_n)\psi'(n+1)^3
  -7168x_n^4-9552x_n^2-\frac{1148}{15}\Bigr\}
  =252\zeta(3).
 \end{aligned}}
\end{equation}
It is again an ordinary absolutely convergent sum. Equations
\eqref{cc:rationalsum} and \eqref{cc:zeta3sum} demonstrate that
the fixed four-coordinate description supplies exact extraction
identities, in addition to individual moment evaluations.

\subsection{The cubic generator from the known square generator}

Let
\[
 \mathcal G_2(t)=\sum_{r\ge0}\mathcal T(-2r)\frac{t^{2r}}{(2r)!},
 \qquad
 \mathcal G_3(t)=\sum_{r\ge0}U_r\frac{t^{2r+1}}{(2r+1)!},
\]
where $\mathcal T(s)=\sum_{n\ge0}A_n^2x_n^{-s}$ for $\Re s>-1$,
continued meromorphically elsewhere, is the centered-square Dirichlet
series of the canonical chapter. The inherited explicit expression is
\begin{equation}\label{cc:G2}
 \mathcal G_2(t)=\frac{t}{2\sinh(t/2)}
 \left[3\zeta(3)-\frac1{4t}
       \int_0^t(t-u)u^2\coth(u/2)\,du\right].
\end{equation}
This is recalled with attribution to \cite{cc:repo-square}, not
presented as an additional new square-moment theorem.

\begin{theorem}[An analytic cubic moment generator]
\label{cc:generator}
The series $\mathcal G_3$ converges for $|t|<2\pi$.  On that disk,
\begin{equation}\label{cc:G3integral}
 \boxed{\displaystyle
 \mathcal G_3(t)=\frac1{2\sinh(t/2)}
 \int_0^t(t-u)\left\{
    3\cosh(u/2)\mathcal G_2(u)-\frac{17}{4}\zeta(4)
    -\frac14\zeta(2)u^2+\frac{u^4}{96}
 \right\}\,du.}
\end{equation}
The expression is holomorphic and odd; its apparent singularity at
zero is removable.  The integration path lies in the disk and is
otherwise arbitrary.
\end{theorem}

\begin{proof}
Put $C(t)=\sum_{j\ge1}C_jt^{2j}/(2j)!$.
The Bernoulli transform \eqref{cc:allformula} gives
\[
 \mathcal G_3(t)=\frac{C(t)}{2\sinh(t/2)}.
\]
Write $E(t)=\tfrac t2\coth(t/2)$ and
\[
 F(t)=\frac1t\int_0^t(t-u)uE(u)\,du
     =\sum_{h\ge0}\frac{(2h+1)B_{2h}}{(2h+3)!}t^{2h+2}.
\]
Direct coefficient multiplication of \eqref{cc:Rj} yields
\begin{align*}
 C''(t)
 &=9\zeta(3)E(t)-\frac32E(t)F(t)
          -\frac{17}{4}\zeta(4)-\frac14\zeta(2)t^2+\frac{t^4}{96}\\
 &=3\cosh(t/2)\mathcal G_2(t)
          -\frac{17}{4}\zeta(4)-\frac14\zeta(2)t^2+\frac{t^4}{96}.
\end{align*}
In the second line use \eqref{cc:G2} and
$E(t)=\cosh(t/2)t/(2\sinh(t/2))$.
All these series converge in $|t|<2\pi$ by the Bernoulli generating
functions.  The initial conditions are $C(0)=C'(0)=0$.
Twice integrating proves \eqref{cc:G3integral} and the stated
analyticity and parity.
\end{proof}

\subsection{A finite depth-two polylogarithmic expression}

For $0<z<1$ use the convention
\[
 \operatorname{Li}_{a,b}(z)
   =\sum_{m>n\ge1}\frac{z^m}{m^a n^b}.
\]
This is the one-variable multiple polylogarithm with the outer index
carrying the argument.  In the following formula, every displayed
depth-two first index is at least two.

\begin{theorem}[Closed polylogarithmic generator]
\label{cc:polylog}
For real $0<t<2\pi$, set $z=e^{-t}$ and define
\begin{align}
 \mathcal C(t)={}&
 \left(9\zeta(4)+6t\zeta(3)-\frac{t^4}{16}\right)
           \operatorname{Li}_2(z)
 +\left(12\zeta(3)-\frac{t^3}{2}\right)\operatorname{Li}_3(z)
       \notag\\
 &-3t^2\operatorname{Li}_4(z)-12t\operatorname{Li}_5(z)
       -\frac{45}{2}\operatorname{Li}_6(z)
       -\frac32t^2\operatorname{Li}_{2,2}(z)
       \notag\\
 &-6t\left(\operatorname{Li}_{3,2}(z)+\operatorname{Li}_{2,3}(z)\right)
       -9\left(\operatorname{Li}_{4,2}(z)+\operatorname{Li}_{2,4}(z)\right)
       -12\operatorname{Li}_{3,3}(z)
       \notag\\
 &+\frac{15}{2}\zeta(6)-6\zeta(3)^2+6t\zeta(5)
       +\frac12t^3\zeta(3)+\frac18t^2\zeta(4)
       -\frac1{48}t^4\zeta(2)-\frac{t^6}{1440}.
       \label{cc:Cpolylog}
\end{align}
Then
\begin{equation}\label{cc:G3polylog}
 \mathcal G_3(t)=\frac{\mathcal C(t)}{2\sinh(t/2)}.
\end{equation}
Equation \eqref{cc:G3integral} specifies its single-valued holomorphic
odd continuation throughout $|t|<2\pi$.
\end{theorem}

\begin{proof}
For ordinary polylogarithms,
$\frac{d}{dt}\operatorname{Li}_a(e^{-t})=-\operatorname{Li}_{a-1}(e^{-t})$.
For the depth-two convention just stated,
\[
 \frac{d}{dt}\operatorname{Li}_{a,b}(e^{-t})
   =-\operatorname{Li}_{a-1,b}(e^{-t})\quad(a\ge2),\qquad
 \frac{d}{dt}\operatorname{Li}_{1,b}(e^{-t})
   =-\frac{\operatorname{Li}_b(e^{-t})}{e^t-1}.
\]
Together with $\coth(t/2)=(1+e^{-t})/(1-e^{-t})$, these give,
by finite differentiation, that $\mathcal C''$ is the integrand in
\eqref{cc:G3integral}.

The integration constants may be checked by anchored primitives.
If $F_a'(t)=-F_{a-1}(t)$ and the displayed endpoint values exist, then
\begin{equation}\label{cc:anchoredprimitive}
 \int_0^t(t-u)u^kF_0(u)\,du
 =\sum_{j=0}^k(j+1)\frac{k!}{(k-j)!}t^{k-j}F_{j+2}(t)
       -(k+1)!F_{k+2}(0)+tk!F_{k+1}(0).
\end{equation}
Apply this with $F_a=\operatorname{Li}_{s+a}(e^{-t})$, or
$F_a=\operatorname{Li}_{a,b}(e^{-t})$ for $k\ge1$.
The only potentially divergent endpoint is kept in the convergent
combination
\[
 \operatorname{Li}_{1,4}(e^{-t})-\zeta(4)\operatorname{Li}_1(e^{-t})
 \longrightarrow -\zeta_{\mathrm{MZ}}(4,1)-\zeta(5)
                 =\zeta(2)\zeta(3)-3\zeta(5).
\]
The elementary depth-two reductions
\begin{align*}
 \zeta_{\mathrm{MZ}}(2,3)&=\tfrac92\zeta(5)-2\zeta(2)\zeta(3),&
 \zeta_{\mathrm{MZ}}(3,2)&=3\zeta(2)\zeta(3)-\tfrac{11}{2}\zeta(5),\\
 \zeta_{\mathrm{MZ}}(4,2)+\zeta_{\mathrm{MZ}}(2,4)&=\zeta(2)\zeta(4)-\zeta(6),&
 2\zeta_{\mathrm{MZ}}(3,3)&=\zeta(3)^2-\zeta(6)
\end{align*}
reduce the endpoint constants to those in \eqref{cc:Cpolylog}.
Thus $\mathcal C(0)=\mathcal C'(0)=0$ in the limiting sense,
and \eqref{cc:G3integral} proves the claim.  No differentiation of
an identity known only at integral zeta orders is used.
\end{proof}

\subsection{Verification scope and immediate continuation questions}

The supplied script independently reconstructs the cutoff constants
using Faulhaber polynomials and the uncentered trigamma expansion; it
does not insert \eqref{cc:Rj} into that reconstruction.  It agrees
with the closed formula for $1\le j\le18$, checks the generator
coefficient conversion over the same range, and verifies the second
derivative of \eqref{cc:Cpolylog} as an exact symbolic zero.
Independent direct tail sums plus centered Hurwitz-tail corrections
check $0\le r\le8$.  The largest displayed residual in the higher
cutoff calculation is below $5\cdot10^{-76}$ at 110-digit working
precision.  Those residuals are numerical diagnostics, not certified
interval bounds or replacements for the finite-sum proof.

The new result suggests three concrete questions.  First, determine
whether all regular centered moments of $\psi'(n+1)^4$ admit a
similarly explicit fixed coordinate space, and derive its generator
with every subtraction and endpoint constant specified.  Second,
identify the analogous regular-parity moment spaces for mixed products
$\prod_i\psi^{(m_i)}(n+1)$, distinguishing formal coefficient rank
from arithmetic independence.  Third, determine whether the particular
depth-two combination in \eqref{cc:Cpolylog} reduces to a prescribed
finite class of classical polylogarithms; a failed numerical search
would not prove that it cannot.  None of these questions is settled
by the cubic calculation above.
